% !TeX root = ../main.tex
\subsection{تمرین \secnum{۴.۶}، اثر بزرگ‌کردن \lr{Inverter}}
\label{sec:delay-exercise-scaled-inverter}

\textbf{خواسته.}
ظرفیت ورودی، \(g\) و \(p\) یک \lr{Inverter} با عرض
ترانزیستورهای چهار برابر را بیابید. ظرفیت ورودی مرجع
\(3C\) و تأخیر بی‌بار نرمال‌شده آن \(p_{\mathrm{inv}}\) است
\cite[تمرین \secnum{۴.۶}، صفحه ۱۷۷]{weste}.

\textbf{حل.}
در مدل تناسب ظرفیت ورودی با عرض، \(C_{\mathrm{in}}\) به
\(4(3C)=12C\) می‌رسد. برای محاسبه \(g\)،
ظرفیت ورودی گیت را به ظرفیت ورودی مرجع تقسیم می‌کنیم.
مرجع برای این مقایسه باید یک \lr{Inverter} هم‌قدرت باشد و در اینجا ظرفیت ورودی آن نیز
\(12C\) است. بنابراین \(g=1\). با فرض تغییر متناسب ظرفیت
\lr{Diffusion}، حاصل‌ضرب مقاومت و ظرفیت داخلی ثابت می‌ماند.
اگر \(C_{\mathrm{diff,inv}}\) ظرفیت داخلی مؤثر
\lr{Inverter} واحد باشد، سهم بی‌بار در تقریب مرتبه اول
چنین محاسبه می‌شود.
\[
  p_{\mathrm{4x}}\approx
  \frac{(R/4)\bigl(4C_{\mathrm{diff,inv}}\bigr)}{\tau}
  =p_{\mathrm{inv}}.
\]
این نتیجه مربوط به بزرگ‌کردن یکنواخت در مدل \lr{RC} است.
با تغییر \lr{Layout}، ظرفیت داخلی می‌تواند از تناسب خطی
با عرض فاصله بگیرد.
